\cotefile{157-5}
\scanpage{20}
\mspage{7}

\diagram{\begin{tikzcd}[ampersand replacement=\&,column sep=large]
 X'_{i_0}\arrow[r]\arrow[d,no head,"\alpha_{i_0}"']\arrow[dr,phantom,"\text{cocart}" description]
 \& X_{i_0}\arrow[r]\arrow[d,"\beta_{i_0}"']\arrow[dr,phantom,"\text{\unsure{cocart}}" description]
 \& X_{\unsure{\infty}}\arrow[d]\\
 Y'_{i_0}\arrow[r,no head] \& \widetilde Y_{i_0}\arrow[r]\arrow[d,"\widetilde u_{i_0}"']\arrow[dr,phantom,"\text{cocart}" description]
 \& Z'_{i_0}\arrow[d,"\unsure{v|u}_{i_0}"]\\
 \& Y_{i_0}\arrow[r,no head] \& Z_{i_0}\,.
\end{tikzcd}}

On a $\widetilde u_{i_0}\in C$ par hypothèse, donc $v_{i_0}\in C$, ce qui achève la démonstration.

\uline{Proposition 2} $(\unsure{\mathcal M|\mathscr M},C)$ et le cas échéant $W$, comme dans Prop.~1 et son corollaire, $I$ ens ordonné localement fini, i.e tel que $\forall i_0\in I$, $I_{\unsure{/|\leq}i_0}$ soit fini. Soit\deleted{\unsure{ent|s}} \deleted{\unsure{y|z}}

\[
 u_\bullet:X_\bullet\longrightarrow Y_\bullet
\]

flèche dans $ {}^{*}\mathcal M(I^\circ)_C={}$\inserted{\unsure{ss cat.} \deleted{\unsure{p}}pleine des} systèmes inductifs sur $I$, \unsure{à valeurs|internes} dans $\mathcal M$, tels
% slanted
\margin{en fait, l’hypothèse que $X_\bullet,Y_\bullet$ dans $\mathcal M(I^\circ)_C$ est inutile, \unsure{tant} dans la prop.~2 que dans les corollaires} que $\varnothing_{\mathcal M(I^\circ)}$\inserted{objet initial de $\mathcal M(I^\circ)$} $\mathbin{\unsure{-|\to}} Z_\bullet$ soit dans $C_I$ (\unsure{défini|définie} dans prop.~1), i.e. tel que $\forall i_0\in I$, la flèche $Z'_{i_0}=\underset{i<i_0}{\underrightarrow{\lim}}\,Z_i\longrightarrow Z_{i_0}$ soit dans $C$, et supposons $u_\bullet\in C_I$. Alors

a) $\underrightarrow{\lim}_{\unsure{I}}u_\bullet:\deleted{\unsure{X}}X_I\overset{\text{déf}}{=}\underset{i\in I}{\underrightarrow{\lim}}\,X_i\longrightarrow Y_I\overset{\text{déf}}{=}\underset{i\in I}{\underrightarrow{\lim}}\,Y_i$ est dans $C$,
% archive sheet number: 19

